%=====================================================================
\chapter{Answers to the Checkpoints}\label{app:answers}
%=====================================================================
\normalfigures

Three answers per chapter, in order. Attempt the checkpoint from memory before
reading these, and argue with them where you disagree --- several are judgement
calls rather than facts, and they say so.

% The slide pipeline parses this appendix: \ansch{<number> --- <title>} must be
% followed immediately by an ans environment of exactly three items.
\newcommand{\ansch}[1]{\subsection*{\color{primary}Chapter #1}}
\newenvironment{ans}
  {\begin{enumerate}[leftmargin=*, itemsep=4pt, topsep=3pt,
                     label=\textbf{\arabic*.}]}
  {\end{enumerate}}

%---------------------------------------------------------------------
\ansch{1 --- What Virtualization Is}
\begin{ans}
  \item Active memory is $24 \times 19 = 456$ GB. With 30\% headroom,
        $456 \times 1.3 = 593$ GB, so $\lceil 593/384 \rceil = 2$ hosts, plus
        one for N+1 $= \mathbf{3}$ hosts. The consolidation ratio is
        $24/3 = \mathbf{8:1}$. Note what was \emph{not} used: the 128 GB
        installed per server. Sizing on that would have given
        $24 \times 128 \times 1.3 / 384 = 11$ hosts for the same work.
  \item The second half is wrong. There \emph{is} a virtualization layer ---
        namespaces and control groups --- it is simply inserted inside the
        kernel rather than beneath it. The correct statement is that containers
        start faster and cost less memory because they share the host's kernel
        instead of each running their own, not because no layer exists.
  \item \textbf{Encapsulation.} A running machine becomes a file, so it can be
        snapshotted, copied, versioned and restored. The operation with no
        physical equivalent is the snapshot: you can return a machine to its
        state four minutes ago in a few seconds. A physical server has no way to
        record its entire state --- memory, disk and devices --- as something
        that can be put back.
\end{ans}

%---------------------------------------------------------------------
\ansch{2 --- Virtualization Architectures}
\begin{ans}
  \item By Goldberg's test it is \textbf{Type-1}: the hypervisor reaches the
        hardware without another operating system's permission, and the stripped
        Linux exists only to boot it. The more useful questions are
        \emph{how large is the trusted computing base?} --- here, a kernel plus
        every driver it loads --- and \emph{what else runs on this machine?}
        Those decide the security and performance arguments; the Roman numeral
        decides neither.
  \item CPU overhead is $(412-400)/400 = \mathbf{3\%}$, which is normal. Disk
        throughput is $1 - 210/1050 = \mathbf{80\%\ lost}$, which is not. The
        pattern --- small CPU penalty, enormous I/O penalty --- is the signature
        of an \textbf{emulated disk controller}. Change the disk bus to virtio;
        expect the loss to fall to a few per cent.
  \item Paravirtualization requires modifying the guest kernel, which requires
        its source and the ability to rebuild it. With the vendor gone, neither
        exists. Use \textbf{hardware-assisted virtualization}, which needs no
        guest changes at all, and install paravirtualized \emph{drivers} if any
        exist for that operating system --- a far weaker requirement than
        recompiling the kernel.
\end{ans}

%---------------------------------------------------------------------
\ansch{3 --- The Virtual Machine Monitor}
\begin{ans}
  \item \texttt{RDCFG} is \textbf{behaviour-sensitive}: its result depends on
        the privilege configuration, which the monitor must control. It is
        \textbf{not privileged}: it does not fault in user mode. So there is a
        sensitive instruction outside the privileged set, the condition
        \emph{sensitive $\subseteq$ privileged} fails, and \textbf{the
        architecture is not virtualizable by trap-and-emulate alone}. A guest
        would read the deprivileged level and could detect, or misbehave in,
        virtualization.
  \item Because a crash is a \emph{signal}. The monitor author would learn
        immediately, at the offending instruction, with a stack trace. What x86
        does instead is accept the instruction, update the arithmetic flags and
        silently discard the write to \texttt{IF}: the guest believes it
        disabled interrupts and did not. The resulting corruption is
        intermittent, unreproducible and traceable to nothing. Silent incorrect
        behaviour is far harder to engineer around than a loud failure.
  \item When you control the guest and want the lowest overhead --- for example
        a cloud provider shipping its own Linux images, or an appliance vendor
        building both halves. A hypercall costs tens of cycles where a trap
        costs hundreds, the monitor is far simpler, and the guest's source is
        available by construction. The requirement that makes paravirtualization
        unusable in general is precisely the one that does not bind here.
\end{ans}

%---------------------------------------------------------------------
\ansch{4 --- Hardware Support for Virtualization}
\begin{ans}
  \item Backward compatibility. Making \texttt{POPF} and the others privileged
        would change the behaviour of existing, correct user programs that rely
        on the write to \texttt{IF} being silently ignored --- they would begin
        faulting. x86's commercial value rests on not breaking old software, so
        both vendors added a new processor \emph{mode} in which existing
        instructions behave as before and the \emph{hardware} decides when to
        hand control to the hypervisor.
  \item For each of the 5 guest levels: 5 EPT accesses to translate that level's
        address, plus 1 read of the guest entry $= 6$, so $5 \times 6 = 30$;
        plus 5 to translate the final guest-physical address:
        $\mathbf{35}$ memory accesses. In general, for $n$ levels:
        $n(n+1) + n = \mathbf{n^2 + 2n}$. Checking against the chapter,
        $n = 4$ gives $16 + 8 = 24$.
  \item \textbf{Device passthrough} is lost --- assigning a physical device
        directly to a guest, including SR-IOV virtual functions. It must be
        refused rather than warned about because a device performs DMA to
        physical addresses: without an IOMMU to translate and confine them, the
        guest's driver can direct the device to write anywhere in host memory,
        including the hypervisor and other guests. That is not a degraded
        configuration; it is the complete loss of the resource-control property
        of \chref{vmm}.
\end{ans}

%---------------------------------------------------------------------
\ansch{5 --- Major Hypervisor Platforms}
\begin{ans}
  \item \textbf{ESXi}: in the VMkernel itself --- a small, purpose-built trusted
        base of about 150 MB, at the cost of a short hardware compatibility
        list. \textbf{Hyper-V}: in the parent partition, a full Windows
        installation --- thin hypervisor, large privileged guest.
        \textbf{KVM}: in the Linux kernel --- every driver Linux has, and all of
        them in the trusted base. \textbf{Xen}: in dom0, or in a separate driver
        domain --- a one-megabyte hypervisor, and the only design that can put a
        driver's blast radius below the whole host.
  \item Docker Desktop enables Hyper-V (or the WSL2 platform, which uses it).
        That slides the Hyper-V hypervisor underneath the running Windows,
        demoting it to a parent partition. VirtualBox then no longer runs on the
        hardware; it runs as a \emph{nested} guest above Hyper-V, and nested
        virtualization multiplies every exit the inner guest takes. Nothing
        about VirtualBox changed --- the machine underneath it did.
  \item Using the chapter's figures, vSphere cost about \$200{,}000 more over
        three years for three hosts, which is roughly 1.5 full-time engineers.
        A licence becomes defensible once the estate is large enough that
        vCenter genuinely saves that much effort --- in practice somewhere in
        the low hundreds of hosts, or sooner where availability has contractual
        penalties. \textbf{The assumption doing the most work} is that a
        self-managed KVM estate would need those 1.5 engineers; if the team
        already runs Linux at scale and has the automation, the saving is much
        smaller and the threshold much higher.
\end{ans}

%---------------------------------------------------------------------
\ansch{6 --- CPU and Memory Virtualization}
\begin{ans}
  \item $\sigma_{\text{cpu}} = (20 \times 6)/48 = 120/48 = \mathbf{2.5}$. A mean
        \%RDY of 14\% is well above the 10\% action threshold, so the host is
        oversubscribed. \textbf{Measure first:} the actual vCPU utilisation of
        each guest --- 6 vCPUs is a large allocation and is very likely more
        than most of them use. \textbf{Then act:} right-size downward before
        buying anything. Halving the allocations would bring $\sigma$ to 1.25
        and would almost certainly resolve the ready time without new hardware.
  \item Because the hypervisor cannot tell which pages matter. From outside,
        every guest page looks alike; it cannot distinguish a hot database
        buffer from a page of cache the guest would discard without noticing,
        so taking pages itself means guessing, and a wrong guess costs the guest
        a disk read. The guest kernel has a page-replacement algorithm that is
        good at exactly this question. The balloon driver asks in the one
        language the guest understands --- it allocates memory --- and the guest
        decides what to give up.
  \item \textbf{The guest is wider than a NUMA node.} It has 40 vCPUs against
        32 cores per socket and 400 GB against 256 GB per socket, so it must
        span both sockets and roughly half its memory accesses cross the
        interconnect at 1.5--2$\times$ the latency. No hypervisor setting fixes
        it because the memory is physically attached to a socket; locality
        cannot be created after the fact. The change that would work is
        \textbf{resizing the guest to fit inside one node} --- 32 vCPU and
        256 GB or less. Failing that, enable vNUMA so the guest's own scheduler
        can at least keep each thread's memory on its own node.
\end{ans}

%---------------------------------------------------------------------
\ansch{7 --- Paravirtualization and Device Virtualization}
\begin{ans}
  \item $1.6/0.2 = \mathbf{8.0}$ host cores per Mpps, roughly twenty-five times
        what vhost costs. That figure, together with the very low packet rate,
        is the signature of an \textbf{emulated network adapter} --- an e1000 or
        similar, where every device register access is a VM exit. Change the
        interface model to \textbf{virtio}, and enable vhost; expect the rate to
        rise by an order of magnitude and the host CPU to fall.
  \item Because the exit is caused by the \emph{doorbell}, not by the request. A
        guest places buffers in the shared virtqueue with ordinary memory
        writes, which cost nothing, and rings the doorbell once to say work is
        waiting --- and one doorbell can cover a hundred queued requests. Under
        load the back-end is already polling the ring and finds new work without
        being told, so the guest can suppress notification entirely. The
        mechanism is \textbf{batching with notification suppression}, and it
        makes the exits-per-request ratio fall as the rate rises.
  \item \textbf{The conflict:} SR-IOV virtual functions are assigned devices, so
        those twelve guests cannot be live-migrated, and live migration is how
        a host is patched without downtime. Two resolutions:
        (i) \emph{bonded failover} --- give each guest a VF and a virtio
        interface in a bond, detach the VF before migrating and reattach it
        afterwards, accepting reduced throughput during the move;
        (ii) \emph{scheduled downtime for those twelve guests only}, documented
        as an exception at design time rather than discovered at the first
        patch window.
\end{ans}

%---------------------------------------------------------------------
\ansch{8 --- From Virtual Machines to Containers}
\begin{ans}
  \item 137 is $128 + 9$: the process was terminated by signal 9, \texttt{SIGKILL},
        which here means the \textbf{memory limit} and the kernel's OOM killer.
        The symptom is a kill rather than a slowdown because memory is not
        divisible in time: a page is either available or it is not, so there is
        nothing to degrade. A CPU limit, by contrast, throttles --- the container
        simply runs more slowly and nothing fails.
  \item A container shares the host's kernel, so there is no guest kernel to
        load a module into; loading one would mean loading it into the
        \emph{host}, which is both a privilege escalation and shared with every
        other container. The two options are (i) load the module on the host and
        expose the resulting device to the container, accepting that the module
        is now host-wide, or (ii) use a \textbf{virtual machine}, or a microVM
        from \chref{lightweight}, which has a kernel of its own.
  \item The image is stored once and shared: $400 + 3 \times 50 =
        \mathbf{550}$ MB. The logs are in each container's writable layer, so
        they count three times and are lost when the containers are replaced.
        The change is to \textbf{write logs to stdout} and let the runtime
        collect them, or mount a volume --- either removes them from the
        copy-on-write layer entirely.
\end{ans}

%---------------------------------------------------------------------
\ansch{9 --- Docker and the Container Ecosystem}
\begin{ans}
  \item \texttt{node:20} is a \emph{tag}, which is a mutable pointer. The
        publisher republished it during the week --- a patch release, a new base
        layer --- so Friday's build used different bytes. The change that
        prevents it is to \textbf{pin the digest}:
        \texttt{FROM node:20@sha256:\ldots}, which identifies exact content for
        ever. Pair it with an automated job that proposes digest updates as
        reviewable changes, so security fixes are still applied deliberately.
  \item \textbf{Same \texttt{RUN}: no.} Only the layer's final state is
        committed, and the key was deleted before the layer closed.
        \textbf{Separate \texttt{RUN}: yes.} The first layer contains the key,
        the second records its deletion, and both layers ship. Anyone who pulls
        the image can extract the earlier layer and read it. This is why secrets
        belong in build secrets or a mount, never in a layer at all.
  \item Stating assumptions: say the dependency layer is 1.0 GB, the application
        layer 20 MB, and there are 20 builds a day.
        \emph{Wrong order:} each build rebuilds and pushes about 1.02 GB, and
        40 nodes pull it: $20 \times 1.02 + 40 \times 20 \times 1.02 \approx
        \mathbf{836}$ GB a day.
        \emph{Right order:} only the 20 MB layer changes:
        $20 \times 0.02 + 40 \times 20 \times 0.02 = \mathbf{16.4}$ GB a day.
        A saving of roughly \textbf{820 GB a day} in transfer alone, before
        counting the build time.
\end{ans}

%---------------------------------------------------------------------
\ansch{10 --- Container Orchestration with Kubernetes}
\begin{ans}
  \item Cluster totals are $6 \times 12 = 72$ cores and $6 \times 48 = 288$ GB;
        after existing requests, 32 cores and 98 GB remain. In aggregate the new
        Deployment's 15 cores and 90 GB would fit --- but Pods are placed on
        individual nodes. Each node has $12 - 40/6 = 5.3$ cores and
        $48 - 190/6 = 16.3$ GB free, giving
        $\min(\lfloor 5.3/1 \rfloor, \lfloor 16.3/6 \rfloor) = 2$ Pods per node.
        So $2 \times 6 = \mathbf{12}$ schedule and 3 stay \texttt{Pending}. The
        binding resource is \textbf{memory}.
  \item \textbf{BestEffort}, which is evicted first whenever a node comes under
        memory pressure. The symptom is that the workloads run perfectly for
        weeks and then are killed --- apparently at random, but in fact
        \emph{only when the cluster is busy}, which is exactly when they are
        most needed. Declaring nothing is not neutral; it is choosing to be
        first in the eviction queue.
  \item A liveness probe \emph{kills} the container when it fails. If the probe
        checks a shared database and the database slows, every replica fails the
        probe at the same moment, so Kubernetes kills them all. They restart,
        reconnect simultaneously in a thundering herd, overload the already-slow
        database and fail again. Use a \textbf{readiness} probe: it removes the
        Pod from the Service's endpoints without killing it, so the replicas
        stop taking traffic, the database recovers, and they return.
\end{ans}

%---------------------------------------------------------------------
\ansch{11 --- Lightweight and Secure Virtualization}
\begin{ans}
  \item Removed: \textbf{firmware initialisation}, \textbf{PCI bus enumeration}
        and \textbf{legacy device emulation}; the kernel is loaded directly into
        memory rather than booted through a bootloader. What the guest must be
        willing to give up is \textbf{looking like a PC} --- it cannot expect a
        BIOS, a PCI bus or any legacy device, so it must be built to use virtio
        and to be entered directly. The cost was compatibility, not
        virtualization.
  \item \textbf{System-call rate.} gVisor handles system calls in a user-space
        Go program instead of in the kernel, so a workload making many small
        calls --- lots of short reads and writes, many file operations --- pays
        heavily, while a CPU-bound workload that rarely calls the kernel is
        barely affected. Confirm by measuring the calls per second of each, with
        \texttt{strace -c} or \texttt{perf}, and correlating that against the
        measured slowdown.
  \item Run the 200 internal services on the default runtime (runc) at full
        density, and the customer build jobs on a hardware-isolated runtime ---
        Kata, or a microVM-backed equivalent. The manifests are otherwise
        identical. The single field that expresses it is
        \textbf{\texttt{runtimeClassName}} in the Pod spec.
\end{ans}

%---------------------------------------------------------------------
\ansch{12 --- Storage Virtualization}
\begin{ans}
  \item Provisioned $= 120 \times 800 = 96$ TB, so
        $\sigma = 96/40 = \mathbf{2.4}$. Written $= 120 \times 240 = 28.8$ TB,
        so headroom $= 40 - 28.8 = \mathbf{11.2}$ TB. Growth is
        $120 \times 20 = 2.4$ TB a month, giving
        $11.2/2.4 = \mathbf{4.7}$ months to exhaustion. Alert at \textbf{80\% of
        the pool} --- 32 TB used, 8 TB free --- because that is reached in about
        six weeks, which is the last point at which there is still time to
        order capacity or reclaim space. Alerting at 95\% leaves no time to do
        anything.
  \item Deleting a file inside a guest only marks blocks free in the
        \emph{guest's} filesystem. The pool underneath was never told, so it
        still holds those blocks as written. The setting is \textbf{discard /
        TRIM, enabled end to end} --- in the guest mount options or a periodic
        \texttt{fstrim}, in the virtual disk controller, and in the storage
        layer. On an estate that has never had it, enabling it commonly returns
        10--20\% of the pool.
  \item \textbf{Slow} because reads are amplified: a block must be looked for in
        each overlay in turn until it is found, so twelve layers can mean twelve
        lookups where one would do, and a first write to a block must copy it up
        from wherever it currently lives. \textbf{Dangerous} because every file
        in the chain is load-bearing: lose or corrupt any one of them and every
        snapshot above it is unreadable --- and after six months nobody knows
        which snapshots are still needed.
\end{ans}

%---------------------------------------------------------------------
\ansch{13 --- Network Virtualization and SDN}
\begin{ans}
  \item $14 + 40 + 8 + (8+16) = \mathbf{86}$ bytes of overhead, so the largest
        inner frame is $1500 - 86 = \mathbf{1414}$ bytes. Note that IPv6 and
        Geneve options together cost 36 bytes more than the VXLAN-over-IPv4 case
        in the chapter --- the arithmetic must be redone for each
        encapsulation, never assumed at 50.
  \item The oversized packet is dropped and an ICMP ``fragmentation needed''
        message is returned to tell the sender to use a smaller size. Path MTU
        discovery depends on that message arriving --- and ICMP is very commonly
        filtered, so the sender never learns, retransmits the same segment and
        stalls for ever. The handshake succeeded because its packets were small.
        The firewall change is to \textbf{permit ICMP type 3 code 4}
        (destination unreachable, fragmentation needed), which turns a silent
        hang into a connection that corrects itself.
  \item A VLAN is a tag within a single layer-2 domain: it does not survive
        being routed, so it cannot span three buildings on different IP subnets.
        It is also limited to 4094 identifiers. Use an \textbf{overlay} ---
        VXLAN or Geneve --- which puts the tenant's Ethernet frame inside an
        ordinary IP packet, so it crosses any routed network, and whose 24-bit
        identifier gives about 16.7 million isolated networks.
\end{ans}

%---------------------------------------------------------------------
\ansch{14 --- Cloud Service and Deployment Models}
\begin{ans}
  \item It satisfies on-demand self-service (departments provision themselves),
        broad network access, resource pooling and measured service (charged per
        guest-month). It fails \textbf{rapid elasticity} --- specifically the
        ability to \emph{shrink}. The hosts are bought; a department returning a
        guest frees a little electricity and nothing else, and the university
        cannot rent the spare capacity to anyone. It is a well-run
        virtualization estate with a cloud-like interface, which is valuable,
        but it does not have the economics.
  \item $10 \times 5 \times 52 = 2{,}600$ hours a year at \$0.2016 is
        $\mathbf{\$524}$ a year. A three-year reservation discounts the rate by
        55\% but commits you to paying continuously:
        $0.45 \times 0.2016 \times 8{,}760 = \$795$ a year. The reservation
        would therefore cost \textbf{more than paying on demand}, because the
        workload uses only $2{,}600/8{,}760 = 30\%$ of the hours it would be
        committing to. Reservations pay for steady load, not for intermittent
        load at any discount.
  \item \textbf{The customer's side} --- under IaaS and under every provider's
        shared responsibility model, data classification and access
        configuration belong to the customer. The provider's default for object
        storage is private; somebody changed it. To stop it recurring: enable
        account-level public-access blocking so it cannot be set at all, scan
        for public buckets continuously, and require infrastructure as code
        (\chref{management}) so that a permission change is a reviewed commit
        rather than a click.
\end{ans}

%---------------------------------------------------------------------
\ansch{15 --- Desktop and Application Virtualization}
\begin{ans}
  \item $70 + 300 \times 3 = \mathbf{970}$ GB. Full clones would be
        $300 \times 70 = 21{,}000$ GB, so the saving is
        $21{,}000/970 \approx \mathbf{21.6:1}$. The base image is read by all
        300 desktops, so it must sit on the fastest storage available and stay
        in cache --- the read amplification of \chref{storage} applies at full
        force.
  \item Steady state is $300 \times 10 = 3{,}000$ IOPS. For the storm: 150
        users over 4 minutes with a 90-second boot gives
        $150 \times (1.5/4) \approx 56$ booting at any instant, so
        $56 \times 120 = 6{,}750$ IOPS, plus the 150 already-logged-in users'
        steady 1{,}500: about $\mathbf{8{,}250}$ IOPS, roughly
        \textbf{2.75 times} the all-day figure --- for four minutes. Size for
        this, not for the average.
  \item \textbf{Latency, not bandwidth.} Every keystroke makes a round trip
        before its character appears, so a high-latency path feels broken no
        matter how much capacity is spare; under about 50 ms it is invisible,
        past 150 ms typing feels detached. Measure the round-trip time from the
        client to the session host, and the protocol's own reported frame
        latency --- not throughput, which the monitoring has already told you is
        fine.
\end{ans}

%---------------------------------------------------------------------
\ansch{16 --- Virtualization Management Platforms}
\begin{ans}
  \item Now: $3 \times 50 + 8 \times 25 = 350$ min a month $= 70$ hours a year.
        After: $3 \times 5 + 8 \times 3 = 39$ min a month $= 7.8$ hours a year,
        plus a 32-hour build. First-year saving
        $= 70 - (7.8 + 32) = \mathbf{30.2}$ hours; thereafter about 62 hours a
        year. \textbf{Adopt} --- but note it is marginal in year one, so the
        case rests on the estate still existing in year two and on the drift and
        review benefits, which are real but not in this arithmetic.
  \item \textbf{In the image:} everything identical across guests --- the
        patched operating system, the agents, the standard configuration.
        \textbf{In user data:} everything that differs --- hostname, network
        address, SSH keys, role. Wrong way round: a hostname baked into the
        image means every clone has the same name and the same SSH host key, so
        they conflict and tooling cannot distinguish them; a security patch
        applied through user data is re-applied at every boot for ever, slowing
        each start and never updating the image everybody else builds from.
  \item \textbf{Convergent} management --- it is the only strategy available for
        a system that cannot be rebuilt. Over the year: write a configuration
        definition that describes it as it currently is, apply it in check-only
        mode until it reports no differences, then bring it under enforcement;
        in parallel, build and test a replacement from that definition, and when
        a rebuild succeeds the machine can move to the immutable regime. The
        intermediate goal is not automation but \emph{a written description that
        is known to be accurate}.
\end{ans}

%---------------------------------------------------------------------
\ansch{17 --- High Availability, Live Migration and Disaster Recovery}
\begin{ans}
  \item The ratio is $400/1{,}200 = \mathbf{0.33}$, so it converges.
        Round 0: $32{,}000/1{,}200 = 26.7$ s, dirtying $10{,}667$ MB.
        Round 1: $10{,}667/1{,}200 = 8.9$ s, dirtying $3{,}556$ MB.
        Round 2: $3{,}556/1{,}200 = 3.0$ s, dirtying $1{,}185$ MB.
        Subsequent rounds are 1.0 s, 0.3 s and 0.1 s, so the total elapsed time
        is about $\mathbf{40}$ seconds with a switchover pause of tens of
        milliseconds.
  \item Each host believes the \emph{other} has failed, because from its own
        position the symptom is identical. \textbf{Without fencing}, both
        restart the other's guests, so the same guest runs twice and both
        instances write to the same disks: not downtime but corruption, which is
        worse than the outage HA was bought to prevent. \textbf{Quorum} requires
        a majority before acting; with two hosts neither side has one, so both
        stop --- which is why a two-node cluster needs a third vote, a witness,
        and why fencing (powering the suspect host off through its management
        controller) must happen before anything is restarted.
  \item $525{,}600 \times 0.0001 = \mathbf{52.6}$ minutes a year. A 90-minute
        planned upgrade \textbf{exceeds the entire annual budget on its own},
        so it is not compatible. Either the upgrade must be performed without
        downtime --- live migration, rolling updates --- or the availability
        target must be renegotiated to 99.9\%, which allows 8 h 46 min.
\end{ans}

%---------------------------------------------------------------------
\ansch{18 --- Performance Monitoring and Capacity Planning}
\begin{ans}
  \item \textbf{\%RDY (or steal time)} and \textbf{balloon / host swap-in}.
        High \%RDY with no memory activity means the guest is not being given a
        processor --- host contention, and the fix is right-sizing or
        rebalancing. Low \%RDY with ballooning or host swap means memory
        pressure, and the fix is reducing overcommit. Both appear from inside
        the guest as ``the application is slow'', and the guest's own 35\% is
        consistent with either.
  \item $1 - 0.95^{12} = 1 - 0.540 = \mathbf{0.46}$, so about \textbf{46\%} of
        page loads contain at least one call slower than the p95. The tail is
        not a rare case for the user; it is close to half of their experience.
  \item With one host down, seven remain, so the figure is
        $76\% \times 8/7 = \mathbf{87\%}$. That is above the 85\% ceiling at
        which reclamation begins, so \textbf{no} --- the cluster should not be
        reported as healthy. It has no real N+1 today, and the 76\% on the
        dashboard conceals exactly that.
\end{ans}

%---------------------------------------------------------------------
\ansch{19 --- Security in Virtualized Environments}
\begin{ans}
  \item Guest compromise: $1/600 = \mathbf{0.17\%}$. VM escape: a host carries
        $600/25 = 24$ guests, so $24/600 = \mathbf{4\%}$. Management plane:
        $\mathbf{100\%}$. \textbf{Defend the management plane first} --- it is
        25 times the reach of an escape, it needs no foothold inside any guest,
        and it is reachable over the network. Do not forget the storage array,
        which holds all 600 guests' disks and is usually less hardened.
  \item Two hardware threads on one core share caches, branch predictors and
        execution resources almost completely, so a side channel between them
        cannot be closed by flushing at VM entry --- the two threads are running
        \emph{simultaneously}, not alternately. Where the neighbours are
        strictly untrusted, the only reliable answer is to stop using the second
        thread. It costs \textbf{15--30\% of the host's apparent capacity}.
        Core scheduling recovers much of that by allowing SMT only between
        siblings of the same guest.
  \item You are being asked to grant \textbf{host root, twice over}.
        \texttt{--privileged} grants all capabilities and disables seccomp and
        AppArmor; the Docker socket allows the agent to start a new container
        with the host filesystem mounted and privileges set, which is the same
        thing by another route. Alternatives: (i) run the agent as a host
        process under systemd with specific, scoped permissions rather than as a
        container; (ii) give it a read-only, scoped API credential --- most
        monitoring needs metrics, not the socket --- and if the vendor cannot
        work without either, run it where it shares a kernel with nothing you
        care about.
\end{ans}

%---------------------------------------------------------------------
\ansch{20 --- GPU and Accelerator Virtualization}
\begin{ans}
  \item The reclamation ladder works because host memory can be recovered
        gradually: pages are shared, then ballooned, then compressed, then
        swapped, and each rung degrades performance rather than causing failure.
        None of that exists for GPU memory. There is no balloon driver, no
        page sharing between contexts and nowhere cheap to swap to, and a
        compute kernel that needs 40 GB of working set cannot run in 20 GB more
        slowly --- the allocation simply fails. GPU memory is a hard allocation,
        not an elastic one.
  \item Demand is $6 \times 5 = 30$ GPU-hours a day against 24 available, so
        \textbf{the card is undersized whatever is done}. \emph{Passthrough}:
        one group at a time, so five groups queue at any moment and the
        utilisation is 100\% with permanently unmet demand. \emph{A six-way
        MIG partition at 12 GB each does not exist} --- profiles are fixed, and
        on an 80 GB card the choices are seven instances of 10 GB (too small for
        12 GB), three of 20 GB, two of 40 GB or one of 80 GB. The realistic
        option is three 20 GB instances, giving three groups concurrent,
        predictable access and a queue for the other three. Predictability is
        better than passthrough; capacity is still short.
  \item \texttt{utilization.gpu} reports the fraction of sampled intervals in
        which \emph{at least one kernel was resident} --- not how much of the
        hardware was doing work. A job using a single streaming multiprocessor
        badly will show 100\% while leaving most of the card idle. Measure
        \textbf{achieved occupancy}, \textbf{memory bandwidth utilisation} and
        the SM activity counters --- through Nsight or the DCGM profiling
        metrics --- which distinguish a busy card from a resident kernel.
\end{ans}

%---------------------------------------------------------------------
\ansch{21 --- Confidential Computing and Emerging Hardware}
\begin{ans}
  \item Because a hostile hypervisor does more than read memory. Encryption
        stops it \emph{reading}, and leaves it able to manipulate the mapping:
        the concrete attack is \textbf{page remapping} --- giving the guest a
        different physical page than it had, or aliasing two guest addresses to
        one physical page --- and \textbf{replay}, re-presenting an old
        encrypted page. The data stays encrypted throughout and the guest is
        compromised anyway. SEV-SNP's reverse map table and TDX's equivalent add
        the integrity check that refuses a contradictory mapping.
  \item They are protected against an operator or an attacker \textbf{reading
        the guest's memory} --- a memory dump, a cold-boot attack, a hypervisor
        inspecting pages. They are \textbf{not} protected against the attack the
        technology exists to stop: an operator who boots a \emph{different}
        image which requests the key at startup. That image's memory is
        perfectly encrypted while it exfiltrates the database. Without
        verifying a measurement, they have bought a performance penalty.
  \item \textbf{Availability} --- the operator can still stop, detain or destroy
        the machine, so a compliance argument resting on continuity is not
        helped at all. \textbf{Side channels} --- timing and cache attacks are
        not fully addressed, and published attacks target these systems
        specifically, so a claim of absolute confidentiality is
        overstated. (Equally acceptable: protection against your own application
        bugs, which leak data the ordinary way, and protection of data once it
        leaves the TEE.)
\end{ans}

%---------------------------------------------------------------------
\ansch{22 --- Case Studies in Enterprise and Cloud Virtualization}
\begin{ans}
  \item At 1.26 the shortfall was small enough to be absorbed by the
        \textbf{upper rungs of the reclamation ladder} --- transparent page
        sharing across 64 similar guests, and ballooning, both of which degrade
        performance imperceptibly. At 1.57 the shortfall exceeded what the guests
        could give back, so the hypervisor reached \textbf{host swapping}, which
        is not a slope but a cliff: guest pages went to disk and latency rose by
        orders of magnitude. The mechanism is the same ladder; the difference is
        which rung it reached.
  \item The cloud \textbf{made existing waste legible} --- on premises an
        over-allocated guest costs nothing once the host is bought, whereas in
        the cloud every unit of waste is a line on a monthly invoice. The two
        items accounting for 75\% were \textbf{lift and shift without
        right-sizing} (44\%) and \textbf{on-demand pricing} (31\%).
  \item Because quorum requires a majority, and two nodes can never produce one:
        a clean split leaves each node with one vote out of two, so neither may
        act and both stop. A third vote at head office breaks the tie. When the
        link to the witness failed, the site \textbf{operated in a documented
        degraded mode with HA disabled} --- accepting that a node failure during
        that window would not be recovered automatically, rather than risking
        the split brain that automatic recovery without quorum would invite.
\end{ans}

%---------------------------------------------------------------------
\ansch{23 --- The Semester Project}
\begin{ans}
  \item Any hypothesis containing a number and capable of being false, for
        example: \emph{``Migrating our six internal services from virtual
        machines to Kubernetes reduces mean deployment time from 25 minutes to
        under 5, without increasing p99 request latency by more than 10\%.''}
        It names two measurable quantities, states thresholds, and could come
        out either way --- which the original could not, because no result would
        have contradicted ``is good for our department''.
  \item That \textbf{the measurement does not support any claim}: the spread of
        the three runs (7\%, 24\%, $-2$\%) is far larger than the 9\% mean
        effect, and one run has the opposite sign. Reporting ``a 9\%
        improvement'' would be reporting noise. The report should say so. Next:
        find the source of the variance --- neighbours, caching, burst credits,
        nesting --- control it, and repeat with enough runs to separate the
        effect from the spread; and if the variance cannot be controlled, report
        that as the finding.
  \item Switch to the captured output shipped in \texttt{code/}. The hypothesis
        is about \emph{partitioning policy and its effect on predictability},
        which can be tested against recorded measurements from a partitioned
        card: the analysis, the percentiles and the comparison with the
        unpartitioned baseline are all still yours to do. State in the method
        section that the figures are from captured runs rather than your own
        hardware, and treat them as comparative. What you must \emph{not} do is
        quietly change the hypothesis to fit what is available.
\end{ans}
